\begin{problem}[第36届全国中学生物理竞赛决赛][汽车轮子破坏性试验]{111}
    新型号汽车在出厂前都要通过破坏性试验。在某次汽车测试中，一个汽车轮子的三根辐条被撞掉了其中一根，轮子的横截面如图所示。该轮子可视为内、外半径分别为
    $R_{1}=\frac{4}{5}R_{0}$ 、 $R_{0}$ 的轮盘和两根辐条组成（假设轮盘可视为匀质环形圆盘，辐条可视为匀质细杆，每根辐条的质量为 $m$。轮盘的质量为 $M=8m$，轮子从图 a 所示位置由静止开始释放，释放时轮子上两辐条所张角的平分线恰好水平，此后该轮子在水平地面上做纯滚动。求
    \begin{subproblem}
        刚释放时轮子的角加速度；
    \end{subproblem}
    \begin{subproblem}
        释放后辐条 OB 首次转到竖直位置时，轮子的角加速度、地面对轮子的摩擦力和支持力。
    \end{subproblem}
\end{problem}